% Generated draft: requires mathematical and source review.
\documentclass[11pt]{article}
\usepackage[margin=25mm]{geometry}
\usepackage{fontspec}
\setmainfont{Latin Modern Roman}
\usepackage{amsmath,amssymb,mathtools,unicode-math}
\setmathfont{Latin Modern Math}
\usepackage{enumitem,needspace,xurl,hyperref}
\setlength{\parindent}{0pt}
\setlength{\parskip}{5pt}
\setlength{\emergencystretch}{3em}
\newcommand{\N}{\mathbb{N}}
\newcommand{\Z}{\mathbb{Z}}
\newcommand{\Q}{\mathbb{Q}}
\newcommand{\R}{\mathbb{R}}
\newcommand{\C}{\mathbb{C}}
\newcommand{\F}{\mathbb{F}}
\newcommand{\A}{\mathbb{A}}
\newcommand{\PP}{\mathbb{P}}
\newcommand{\E}{\mathbb{E}}
\newcommand{\eps}{\varepsilon}
\newcommand{\dd}{\,\mathrm{d}}
\begin{document}
\section*{UnsolvedMath ID: 30004119}
\textbf{Source problem code:} catalogue entry
\subsection{Definitions and mathematical statement}
A packing of unit disks in the Euclidean plane is a family of closed disks \(\{\overline B(x,1):x\in X\}\), where \(X\subset\R^2\) and \(\lVert x-y\rVert\geq 2\) for distinct \(x,y\in X\); thus disk interiors are disjoint. For a parameter \(\varepsilon>0\), define the undirected graph
\[
G_\varepsilon(X)=(X,E_\varepsilon),\qquad
E_\varepsilon=\bigl\{\{x,y\}:x,y\in X,\ x\ne y,\ \lVert x-y\rVert\leq 2+\varepsilon\bigr\}.
\]
An infinite connected component means a connected component containing infinitely many centers. The triangular-lattice packing of unit disks has covered-area density \(\rho_\triangle=\pi/\sqrt{12}\).

The supplied question asks for quantitative density assumptions, expressing proximity to \(\rho_\triangle\), that force \(G_\varepsilon(X)\) to have an infinite connected component. A threshold formulation to be clarified is the following: for a specified packing class \(\mathcal P\), a specified covered-area density functional \(D\), and a specified \(\varepsilon>0\), does there exist \(\delta>0\) such that
\[
\forall X\in\mathcal P,\qquad
D(X)>\rho_\triangle-\delta
\ \Longrightarrow\
G_\varepsilon(X)\text{ has an infinite connected component}?
\]
If such a formulation is intended, determine sufficient quantitative bounds on \(\delta\), including their dependence on \(\varepsilon\) and any model parameters.

This is an unresolved formulation schema, not an assertion that such a threshold exists. The record does not specify the density functional, the admissible packing class, or the quantification over \(\varepsilon\). Its statements concern packings without an explicit probability model, whereas its literature assessment identifies a conjecture about random packings and refers to Gibbs states. If a probabilistic formulation is intended, the probability law and the required probability of percolation must also be specified. No deterministic or almost-sure conclusion is silently substituted for the other.
\subsection{Short English statement}
Join distinct centers of nonoverlapping unit disks whenever their distance is at most \(2+\varepsilon\). The question asks how close the packing density must be to the triangular-lattice value to force an infinite connected component. A precise threshold question requires a density definition, a class of packings or a probability model, and quantifiers for \(\varepsilon\); these are missing from the supplied record.
\subsection{Sources}
\begin{itemize}
\item Local UnsolvedMath registry, permanent ID \(30004119\), title: Percolation Structure in Nearly Optimal Disk Packings; exact source problem code: \texttt{catalogue entry}. The record belongs to the Oberwolfach Reports Open Problems collection and gives proposed year 2019. No individual proposer is supplied.
\item Supplied source citation: Mini-Workshop: Mathematics of Crystallisation (2020), Oberwolfach Reports, DOI 10.4171/owr/2019/18, pp. 1071--1111. \url{https://doi.org/10.4171/owr/2019/18}. The literature-source label for this same link is ``Oberwolfach Report 18/2019, Percolation of two-dimensional hard disks.'' The supplied 2020 workshop label and 2019 report designation are retained without resolving their discrepancy.
\item Supplied literature source: Bowen--Lyons--Radin--Winkler, A solidification phenomenon in random packings. \url{https://doi.org/10.1137/050647785}. The registry attributes the conjectural identification to these authors in its literature assessment; it does not identify an individual proposer of the registry problem.
\item Literature assessment date: \texttt{2026-08-22}. This is supplied dated curation, not independently verified source consultation or a current literature determination. The supplied assessment labels the problem open, identifies it with the Bowen--Lyons--Radin--Winkler percolation conjecture, and says that available density and Gibbs-state results support it but do not prove the assertion for every near-optimal packing model. The background explicitly requests independent verification of time-sensitive status and cited claims.
\end{itemize}
\end{document}
